\section{Operationalization of Cognitive and Transfer Outcomes}

\subsection{Matched task families}
Three candidate task families operationalize executive control across the stationary (S), embodied (E), and embodied augmented-reality (A) contexts. First, a go/no-go rule requires responding to designated target cues and withholding responses to distractors. Second, a rule-switching task requires selecting between two explicitly cued stimulus-response mappings. Third, a sequence-updating task requires retaining and revising a short target sequence. The final battery, age-appropriate difficulty, stimulus timing, and number of trials must be established in pilot work rather than selected after observing treatment effects.

Each family uses identical logical rules and target frequencies across contexts. In S, a touchscreen response encodes the decision. In E, a controlled phone orientation or movement encodes it without environmental augmentation. In A, the corresponding action targets a spatially anchored cue. These interfaces are not assumed to be psychometrically equivalent: differences in response latency, reach distance, visual search, and motor precision must be quantified before comparisons of executive performance.

\subsection{Trial-level event schema}
Each trial should record a pseudonymous participant and session identifier, assigned condition, task family, rule identifier, challenge level, cue onset, response event, response correctness, omission or commission error, controller state, movement summary, tracking-quality flag, and safety interruption. All event timestamps must use a documented monotonic device clock; wall-clock time is not a substitute for within-trial timing. Missing responses, timeouts, and invalid sensor events receive separate codes. Environmental video and raw camera frames are not required for these analyses.

\subsection{Outcome definitions}
For go/no-go tasks, commission-error rate is the number of responses on no-go trials divided by the number of valid no-go opportunities; omission-error rate is the number of absent responses on go trials divided by valid go opportunities. For switching tasks, the switch cost is the difference in response latency between valid switch and repeat trials under prespecified accuracy restrictions. Sequence updating can be summarized by the proportion of fully correct sequences, with length-specific denominators. Speed and accuracy are reported together to avoid interpreting impulsive rapid responses as improved control.

Movement exposure is quantified separately from cognitive accuracy. Candidate measures include cumulative angular displacement, movement duration, and timing error relative to the task cue. Tracking loss and frame-rate degradation are fidelity outcomes, not evidence of poor executive function. Where motor execution differs materially across interfaces, cognitive outcomes require an explicit task-equivalence gate before causal interpretation.

\subsection{Independent transfer battery}
The transfer assessment must use an untrained stimulus set and response format that does not reproduce the AR spatial search or phone movement. A candidate battery includes a computer- or examiner-administered inhibition task and an independently presented rule-switching task. Their age appropriateness, test-retest reliability, practice effects, administration burden, and scoring must be evaluated before design freeze. An assessment using a different display but the same practiced stimulus-response mapping does not meet the independence requirement.

The primary endpoint must be one prespecified score, not whichever task yields the largest improvement. Candidate choices include a validated age-standardized composite or a single reliability-qualified inhibition outcome. Secondary transfer outcomes can test rule switching and retention, with a multiplicity policy fixed prospectively. No named instrument is asserted to be validated for this protocol without a source-specific psychometric review.

\subsection{Quality control and interpretability}
Pilot acceptance criteria should address comprehension, ceiling and floor effects, usable trial proportion, response-device latency, and cross-context challenge equivalence. Thresholds are to be set before the confirmatory experiment. Report denominators at the child, session, block, and trial levels. A higher movement score without preserved cognitive performance is not an executive-control benefit; improved trained-task performance without an independent endpoint effect is not evidence of transfer.